\documentclass[sigconf]{acmart}
%%
%% \BibTeX command to typeset BibTeX logo in the docs
\AtBeginDocument{%
  \providecommand\BibTeX{{%
    Bib\TeX}}}

%% Rights management information.  This information is sent to you
%% when you complete the rights form.  These commands have SAMPLE
%% values in them; it is your responsibility as an author to replace
%% the commands and values with those provided to you when you
%% complete the rights form.
\setcopyright{acmlicensed}
\copyrightyear{2018}
\acmYear{2018}
\acmDOI{XXXXXXX.XXXXXXX}

%% These commands are for a JOURNAL article.
\acmJournal{JDS}
\acmVolume{37}
\acmNumber{4}
\acmArticle{111}
\acmMonth{8}

%%
%% Submission ID.
%% Use this when submitting an article to a sponsored event. You'll
%% receive a unique submission ID from the organizers
%% of the event, and this ID should be used as the parameter to this command.
%%\acmSubmissionID{123-A56-BU3}



%%
%% end of the preamble, start of the body of the document source.
\begin{document}

%%
%% The "title" command has an optional parameter,
%% allowing the author to define a "short title" to be used in page headers.
\title{A Survey of Student Sentiments on AI-Assisted Software Engineering Courses}

%%
%% The "author" command and its associated commands are used to define
%% the authors and their affiliations.
%% Of note is the shared affiliation of the first two authors, and the
%% "authornote" and "authornotemark" commands
%% used to denote shared contribution to the research.

\author{Skyler Oakeson}
\affiliation{%
	\institution{Utah State University}
	\city{Logan}
	\state{UT}
	\country{USA}}
\email{skyler.oakeson@usu.edu}

\author{Seth Poulsen}
\affiliation{%
	\institution{Utah State University}
	\city{Logan}
	\state{UT}
	\country{USA}}
\email{seth.poulsen@usu.edu}

%%
%% By default, the full list of authors will be used in the page
%% headers. Often, this list is too long, and will overlap
%% other information printed in the page headers. This command allows
%% the author to define a more concise list
%% of authors' names for this purpose.
\renewcommand{\shortauthors}{Oakeson et al.}
%%
%% Article type: Research, Review, Discussion, Invited or Position
\acmArticleType{Research}

%% Authors' contribution
\acmContributions{BT and GKMT designed the study; LT, VB, and AP
	conducted the experiments, BR, HC, CP and JS analyzed the results,
	JPK developed analytical predictions, all authors participated in
	writing the manuscript.}

\keywords{Do, Not, Use, This, Code, Put, the, Correct, Terms, for,
	Your, Paper}

\begin{abstract}
	This is a paper
\end{abstract}

\maketitle


\section{Introduction}

In this document we discuss how to write an ACM article.

\section{Related Work}

Research on generative AI (GenAI) in programming and software engineering
education has grown quickly, and can be organized into three broad strands that
motivate the present survey: (1) how courses and curricula are being redesigned
around AI-assisted development, (2) how students perceive, use, and feel about
AI tools while learning to program, and (3) how these perceptions relate to
measurable outcomes such as performance, self-efficacy, and skill development.

\subsection{Redesigning Courses and Curricula for AI-Assisted Development}

As generative AI coding tools have become embedded in professional software
development, universities have begun to design dedicated courses and to rework
existing ones around AI-assisted workflows. Geng et al.~\cite{geng2026mapping}
analyzed the syllabi and course materials of 23 U.S. upper-division,
credit-bearing courses that explicitly frame generative AI as part of software
engineering, characterizing their learning objectives, assessment methods,
topics, and the AI tools they document, in order to map how this emerging
curricular area is currently being defined. Complementing this curriculum-level
view, B\"ottcher et al.~\cite{bottcher2025concepts} argue that because students
already use generative AI tools regardless of instructor attitudes toward them,
educators need to explicitly incorporate systematic and professional AI-tool
use into teaching and learning, which requires rethinking constructive
alignment---that is, the fit between learning objectives, teaching and learning
activities, and assessments---for introductory software development courses.
Related work in this space has examined how instructors integrate AI tools into
programming and software engineering curricula more broadly, including concerns
about maintaining core competencies such as debugging, algorithmic thinking,
and independent problem-solving alongside AI-assisted
instruction~\cite{kohenvacs2025integrating}.

\subsection{Student Perceptions and Patterns of AI Use in Programming Courses}

A second strand of work surveys or otherwise elicits how students themselves
perceive and use AI tools once they are integrated into programming
instruction. Kohen-Vacs, Usher, and Jansen~\cite{kohenvacs2025integrating}
surveyed undergraduate students on the usefulness of GenAI tools for
coursework, their potential to enhance creativity, students' behavioral
intention to continue using them, and associated concerns; students reported
generally favorable, medium-to-high perceptions of usefulness and creativity
support and strong intentions to keep using GenAI tools, while expressing
comparatively modest concerns centered mainly on the risk of over-reliance. In
a complementary study conducted within introductory programming labs, G\"uner
and Er~\cite{guner2025classroom} identified five distinct profiles of
student--AI interaction---ranging from AI-reliant code generation to more
independent or collaborative styles---and showed that these profiles shifted
substantially once students received explicit training in prompting and were
given example prompts, with better-performing profiles becoming more common as
guidance increased. This finding echoes a broader theme in the literature that
unguided or unstructured use of GenAI tools can be less beneficial than guided,
scaffolded use, a pattern also reported for structured prompt training in
engineering and data-analysis courses.

ther work has probed student sentiment through more targeted or qualitative
instruments. P\u{a}durean et al.~\cite{padurean2026understanding}, studying
more than 900 CS1 students attempting dialogue-based ``Prompt
Problems''---exercises in which students write natural-language prompts for a
code-generating model rather than code itself---found that students generally
rated prompting as easier and more enjoyable than traditional coding tasks and
perceived it as well suited to developing problem-solving skills, while also
cataloguing the specific prompt-level mistakes (chiefly omission of key
details) that led to unsuccessful attempts. Thorgeirsson, Weidmann, and
Su~\cite{thorgeirsson2026computer} took a related but distinct angle, running
a controlled, preregistered study of ``vibe coding''---programming purely
through natural-language prompting and observed behavior, without viewing the
underlying code---and found that both written-communication skill and
computer-science achievement independently and significantly predicted
vibe-coding performance, with prompt quality mediating much of the association
between writing skill and outcomes.

\subsection{Linking Student Attitudes to Performance and Self-Efficacy}

A third strand connects student sentiment more directly to measured outcomes.
Prather et al.~\cite{prather2026validated} developed and validated a
self-efficacy scale for programming with generative AI, adapting the
established Steinhorst self-efficacy subscales (covering tracing, patterns,
problem solving, and control) to a GenAI-forward introductory computing
context, and found that these foundational constructs remained reliable even
when the surrounding curriculum incorporated AI tools throughout. Kohen-Vacs,
Usher, and Jansen~\cite{kohenvacs2025integrating} additionally examined
students' technical performance on a task requiring them to correct flawed,
LLM-generated code during exams, finding that students performed significantly
worse on these debugging tasks than on comparable instructor-designed
programming and theoretical exam components---a gap that surfaced despite
students' generally positive stated attitudes toward GenAI tools, underscoring
a divergence between favorable sentiment and demonstrated competence in
critically evaluating AI output. Taken together with the profile- and
performance-linked findings of G\"uner and Er~\cite{guner2025classroom}, this
body of work suggests that student sentiment toward AI-assisted instruction
cannot be assumed to track straightforwardly with learning gains, motivating
survey instruments---like the one used in the present study---that probe
sentiment as a construct distinct from, but related to, performance and
self-efficacy.

\subsection{Positioning of the Present Study}

Existing research thus offers converging evidence that (a) AI-assisted software
development courses are a newly forming and heterogeneous curricular category,
(b) students hold generally favorable but multidimensional attitudes toward AI
tools that vary with prior instruction, guidance, and prompting skill, and (c)
these attitudes are only loosely coupled with objective performance on
AI-related tasks. However, much of this work either surveys perceptions in a
single course or intervention, or infers sentiment indirectly from interaction
logs and performance data rather than through validated, retrospective
self-report instruments aimed specifically at students who have completed a
full AI-assisted software development course. The present study addresses this
gap by directly surveying students who have been enrolled in such courses about
their sentiments toward the experience, complementing the curriculum-mapping,
interaction-profiling, and performance-linked perspectives established in prior
work.


\section{Methods}

\section{Results}

\section{Discussion}

\section{Conclusion}


\bibliographystyle{ACM-Reference-Format}
\bibliography{references}

\end{document}
\endinput
